\documentclass[11pt,a4paper]{article}
\usepackage[T1]{fontenc}
\usepackage[utf8]{inputenc}
\usepackage{lmodern,microtype,amsmath,amssymb,amsthm,mathtools}
\usepackage[margin=23mm,headheight=15pt]{geometry}
\usepackage{booktabs,array,longtable,fancyhdr}
\usepackage[hyphens]{url}
\usepackage[hidelinks]{hyperref}
\newtheorem{theorem}{Theorem}[section]
\newtheorem{lemma}[theorem]{Lemma}
\newtheorem{proposition}[theorem]{Proposition}
\newtheorem{corollary}[theorem]{Corollary}
\theoremstyle{remark}\newtheorem{remark}[theorem]{Remark}
\numberwithin{equation}{section}
\allowdisplaybreaks
\setlength{\emergencystretch}{3em}
\setcounter{tocdepth}{1}
\pagestyle{fancy}\fancyhf{}
\fancyhead[L]{\small Further Sun conjectures}
\fancyhead[R]{\small Complete analytic proofs}
\fancyfoot[C]{\thepage}
\hypersetup{pdftitle={Further Sun Conjectures: Complete Proofs},pdfauthor={Keio Research Programme}}
\title{Further Sun Conjectures: Complete Proofs\\[5pt]
\large Parameter families and exact algebraic certificates}
\author{Keio Research Programme\\Research derivation and independent internal verification}
\date{19 November 2026}
\begin{filecontents*}[overwrite]{verify_beta_specializations.py}
#!/usr/bin/env python3
"""Exact algebraic certificates for Sun 4.5, 4.7, 4.9, 4.11.

The proof is in proof_notes.tex. SymPy verifies polynomial identities and
rational functions, not finite approximations to the claimed proof.
The final mpmath check is a separate numerical transcription safeguard.
"""
import json
from pathlib import Path
import sympy as S
import mpmath as mp

w,t,k,a=S.symbols('w t k a')
s=w*w
D=lambda f:S.factor(-w*(1-w*w)/(2*(1+2*w*w))*S.diff(f,w))
B=w/(1+2*w*w)
L=-(1-w*w)/(2*(1+2*w*w))
R=S.Rational
rows=[
 dict(conjecture='4.5', equations=('4.9','4.10'), x=R(1,24), r=R(-1,2), u=S.Integer(3), lam=3,
      c=(3,0,1), P=(7,10,3), Q=(-2,-4), K=3*s*s+1,
      A=(49,-146,21), C=(-3038,1160),
      U=w*(3*s*s+1)*(1292*s*s-16*s+425)/(4*(1+2*s)**5),
      V=54*w*(99*s+46)*(3*s*s+1)/(1+2*s)**4,
      endpoint_U=5*S.sqrt(3),endpoint_V=216*S.sqrt(3)),
 dict(conjecture='4.7', equations=('4.11','4.12'), x=R(-1,25), r=R(5,4), u=R(1,5), lam=5,
      c=(5,14,25), P=(133,131,26), Q=(76,40), K=5*s*s-14*s+25,
      A=(21413,-1409,1036), C=(R(472948,23),R(69280,23)),
      U=w*(5*s*s-14*s+25)*(34900*s*s-15056*s+1135)/(4*(1+2*s)**5),
      V=18*w*(8033*s-1538)*(5*s*s-14*s+25)/(23*(1+2*s)**4),
      endpoint_U=-100*S.sqrt(5),endpoint_V=R(1440,23)*S.sqrt(5)),
 dict(conjecture='4.9', equations=('4.13','4.14'), x=R(-1,72), r=R(3,2), u=R(1,3), lam=3,
      c=(3,8,9), P=(55,54,11), Q=(22,12), K=3*s*s-8*s+9,
      A=(3575,-1026,67), C=(R(42350,13),R(2904,13)),
      U=w*(3*s*s-8*s+9)*(11668*s*s-6896*s+655)/(4*(1+2*s)**5),
      V=18*w*(1571*s-482)*(3*s*s-8*s+9)/(13*(1+2*s)**4),
      endpoint_U=-15*S.sqrt(3),endpoint_V=R(216,13)*S.sqrt(3)),
 dict(conjecture='4.11', equations=('4.16','4.15'), x=R(1,128), r=S.Integer(-1), u=S.Integer(2), lam=2,
      c=(8,8,4), P=(40,44,11), Q=(-8,-8), K=2*s*s-2*s+1,
      A=(200,76,-17), C=(-5800,392),
      U=-w*(2*s*s-2*s+1)*(188*s*s+1184*s+5)/(1+2*s)**5,
      V=72*w*(103*s+44)*(2*s*s-2*s+1)/(1+2*s)**4,
      endpoint_U=-5*S.sqrt(2),endpoint_V=144*S.sqrt(2))]

def iszero(f):
    assert S.cancel(f)==0, S.factor(f)

# Check the universal first-order differential substitutions exactly.
x_w=-(1-w*w)**3/(64*w*w)
iszero(D(x_w)-x_w)
iszero(D(S.log(w))-L)
for n in range(12):
    ck=64**n*S.rf(a,2*n)*S.rf(1-a,n)/S.factorial(3*n)
    h=2*S.harmonic(4*n)-3*S.harmonic(2*n)+S.harmonic(n)
    iszero(ck.subs(a,R(1,2))-S.binomial(4*n,n))
    iszero(S.diff(ck,a).subs(a,R(1,2))-S.binomial(4*n,n)*h)

mp.mp.dps=100
results=[]
for row in rows:
    x,r,u=row['x'],row['r'],row['u']
    c0,c1,c2=row['c']
    assert 0<abs(x)<R(27,256)
    W=c0*(1-t)**2+c1*t*(1-t)+c2*t*t
    iszero(1-64*x*t*t*(1-t)-W/c0*(1-t/r))
    iszero(u-(1-1/r))
    pc=c0*(k+R(1,2))*(k+R(3,2))+c1*(2*k+R(1,2))*(k+R(1,2))+c2*(2*k+R(1,2))*(2*k+R(3,2))
    qc=c0*(-2*k-2)-c1*k+c2*(4*k+2)
    p2,p1,p0=row['P'];q1,q0=row['Q']
    iszero(pc-(p2*k*k+p1*k+p0));iszero(qc-(q1*k+q0))
    iszero((1-w*w)**3+64*x*w*w+(s-u)*row['K']/row['lam'])
    z=S.symbols('z')
    assert S.discriminant(row['K'].subs(w*w,z),z)<0
    aa,bb,cc=row['A'];dd,ee=row['C']
    U=aa*D(D(B))+bb*D(B)+cc*B
    V=-4*aa*D(B)*L-2*aa*B*D(L)-2*bb*B*L+dd*D(B)+ee*B
    iszero(U-row['U']);iszero(V-row['V'])
    iszero(U.subs(w,S.sqrt(u))-row['endpoint_U'])
    iszero(V.subs(w,S.sqrt(u))-row['endpoint_V'])
    # 400 terms at 100 decimal places: numerical safeguard only.
    xx=mp.mpf(str(x.p))/int(x.q)
    hnum=mp.mpf(0);sum_den=mp.mpf(0);sum_poly=mp.mpf(0)
    bkx=mp.mpf(1)
    harmonics=[mp.mpf(0)]
    for n in range(1,1601):harmonics.append(harmonics[-1]+mp.mpf(1)/n)
    for n in range(400):
        hn=2*harmonics[4*n]-3*harmonics[2*n]+harmonics[n]
        sum_den += bkx*((p2*n*n+p1*n+p0)*hn+q1*n+q0)/((3*n+1)*(3*n+2))
        df=mp.mpf(str(S.numer(dd)))/int(S.denom(dd))
        ef=mp.mpf(str(S.numer(ee)))/int(S.denom(ee))
        sum_poly+=bkx*((aa*n*n+bb*n+cc)*hn+df*n+ef)
        bkx*=xx*mp.mpf((4*n+1)*(4*n+2)*(4*n+3)*(4*n+4))/((n+1)*(3*n+1)*(3*n+2)*(3*n+3))
    uu=mp.mpf(str(u.p))/int(u.q)
    rhs_den=-c0*mp.log(uu)/mp.sqrt(uu)
    eu=mp.mpf(str(S.N(row['endpoint_U'],105)))
    ev=mp.mpf(str(S.N(row['endpoint_V'],105)))
    rhs_poly=-eu*mp.log(uu)+ev
    err_den=abs(sum_den-rhs_den);err_poly=abs(sum_poly-rhs_poly)
    assert err_den<mp.mpf('1e-90')
    assert err_poly<mp.mpf('1e-90')
    result={'conjecture':row['conjecture'],'equations':row['equations'],
            'exact_certificates':'PASS',
            'numerical_denominator_residual':mp.nstr(err_den,5),
            'numerical_polynomial_residual':mp.nstr(err_poly,5),
            'denominator_value':mp.nstr(rhs_den,40),
            'polynomial_value':mp.nstr(rhs_poly,40)}
    results.append(result)
    print(json.dumps(result))
out=Path(__file__).with_name('verification_results.json')
out.write_text(json.dumps({'all_exact_certificates_passed':True,'results':results},indent=2)+'\n')
print('ALL EXACT CERTIFICATES AND EIGHT NUMERICAL TRANSCRIPTION CHECKS PASSED')
\end{filecontents*}
\begin{filecontents*}[overwrite]{verify_beta_general_families.py}
"""Exact symbolic-u audit of the two universal beta families."""
import sympy as S

u, w, t, k, z = S.symbols("u w t k z")
s = w*w


def op(f):
    return S.cancel(-w*(1-s)/(2*(1+2*s))*S.diff(f, w))


B = w/(1+2*s)
L = -(1-s)/(2*(1+2*s))
DB, DDB, DL = op(B), op(op(B)), op(L)

A = 4*(u-4)*(2*u-5)**2*(2*u+1)**2
b = 2*(2*u-5)*(12*u**4-76*u**3+237*u**2-267*u+40)
c = (u-1)**2*(2*u-5)*(4*u**2-24*u+15)
d = -4*(u-4)*(u-1)*(2*u+25)*(2*u+1)**2
e = -8*(u-1)**2*(28*u**2-44*u+25)
K = u*s*s+u*(u-3)*s+1
R = -(2*u-5)*(
    4*(20*u*u+176*u-385)*s*s
    -16*(34*u*u-95*u-20)*s
    +5*(28*u*u-56*u+1)
)
T = -72*((8*u**3+20*u*u-58*u+75)*s-2*u*(4*u*u-32*u+37))
U = S.cancel(A*DDB+b*DB+c*B)
V = S.cancel(-4*A*DB*L-2*A*B*DL-2*b*B*L+d*DB+e*B)
assert S.cancel(U-w*K*R/(1+2*s)**5) == 0
assert S.cancel(V-w*K*T/(1+2*s)**4) == 0
print("PASS symbolic-u U and V factor certificates")

# First cancel the common w factor, then evaluate the rational function in w^2.
assert S.cancel(S.cancel(U/w).subs(w*w, u)+5*(u-1)**2*(2*u-5)) == 0
assert S.cancel(S.cancel(V/w).subs(w*w, u)+72*u*(u-1)**2) == 0
print("PASS symbolic-u surviving-root values")

x = (u-1)**3/(64*u)
assert S.cancel((1-s)**3+64*x*s+(s-u)*K/u) == 0
assert S.factor(S.discriminant(u*z*z+u*(u-3)*z+1, z)-u*(u-1)**2*(u-4)) == 0
assert S.factor(S.diff((u-1)**3/u, u)-(u-1)**2*(2*u+1)/u**2) == 0
print("PASS root factorization, discriminant, and domain monotonicity")

W = 1-(u-1)*t+(u-1)**2/u*t*t
assert S.cancel(W-((1-t)**2+(3-u)*t*(1-t)+t*t/u)) == 0
assert S.cancel(1-64*x*t*t*(1-t)-(1+(u-1)*t)*W) == 0
P = (
    (k+S.Rational(1, 2))*(k+S.Rational(3, 2))
    +(3-u)*(2*k+S.Rational(1, 2))*(k+S.Rational(1, 2))
    +(2*k+S.Rational(1, 2))*(2*k+S.Rational(3, 2))/u
)
R0 = (-2*k-2)-(3-u)*k+(4*k+2)/u
P0 = (
    (7-2*u+4/u)*k*k
    +(S.Rational(13, 2)-3*u/2+4/u)*k
    +S.Rational(3, 2)-u/4+S.Rational(3, 4)/u
)
assert S.cancel(P-P0) == 0
assert S.cancel(R0-((u-5+4/u)*k-2+2/u)) == 0
print("PASS denominator-family W/Q factorization and P/R coefficients")

for u0, scale, row in [
    (3, -4, [49, -146, 21, -3038, 1160]),
    (S.Rational(1, 5), S.Rational(-92, 3125),
     [21413, -1409, 1036, S.Rational(472948, 23), S.Rational(69280, 23)]),
    (S.Rational(1, 3), S.Rational(-52, 243),
     [3575, -1026, 67, S.Rational(42350, 13), S.Rational(2904, 13)]),
    (2, -1, [200, 76, -17, -5800, 392]),
]:
    assert all(S.cancel(p.subs(u, u0)-scale*q) == 0 for p, q in zip([A,b,c,d,e], row))
    print(f"PASS five-vector scaling u={u0}, scalar={scale}")

assert c.subs(u, 1) == d.subs(u, 1) == e.subs(u, 1) == 0
assert R0.subs(u, 1) == 0
assert all(p.subs(u, S.Rational(5, 2)) == 0 for p in [A, b, c])
print("PASS exceptional parameter checks u=1 and u=5/2")
print("All independent general-family checks passed.")
\end{filecontents*}
\begin{filecontents*}[overwrite]{verify_beta_independent.py}
"""Independent exact checks for the eight Sun II beta-integral closures.

This checker reconstructs the differential expressions from their definitions.
It does not import the proof author's computed factorizations. Analytic domain,
endpoint domination, and logarithm branches are reviewed in proof_audit.md.
"""

import sympy as S

w, t, k = S.symbols("w t k")
s = w**2


def op(f):
    return S.cancel(-w * (1 - s) / (2 * (1 + 2 * s)) * S.diff(f, w))


B = w / (1 + 2 * s)
L = -(1 - s) / (2 * (1 + 2 * s))
DB, DDB, DL = op(B), op(op(B)), op(L)

rows = [
    (
        S.Rational(1, 24), S.Integer(3), (3, 0, 1),
        (7*k*k + 10*k + 3, -2*k - 4),
        (49, -146, 21, -3038, 1160),
        3*s*s + 1, 5*S.sqrt(3), 216*S.sqrt(3),
    ),
    (
        S.Rational(-1, 25), S.Rational(1, 5), (5, 14, 25),
        (133*k*k + 131*k + 26, 76*k + 40),
        (21413, -1409, 1036, S.Rational(472948, 23), S.Rational(69280, 23)),
        5*s*s - 14*s + 25, -100*S.sqrt(5), S.Rational(1440, 23)*S.sqrt(5),
    ),
    (
        S.Rational(-1, 72), S.Rational(1, 3), (3, 8, 9),
        (55*k*k + 54*k + 11, 22*k + 12),
        (3575, -1026, 67, S.Rational(42350, 13), S.Rational(2904, 13)),
        3*s*s - 8*s + 9, -15*S.sqrt(3), S.Rational(216, 13)*S.sqrt(3),
    ),
    (
        S.Rational(1, 128), S.Integer(2), (8, 8, 4),
        (40*k*k + 44*k + 11, -8*k - 8),
        (200, 76, -17, -5800, 392),
        2*s*s - 2*s + 1, -5*S.sqrt(2), 144*S.sqrt(2),
    ),
]

for x, u, cs, target_pr, weights, K, expected_u, expected_v in rows:
    assert abs(x) < S.Rational(27, 256)
    c0, c1, c2 = cs
    a, b, c, d, e = weights
    W = c0*(1-t)**2 + c1*t*(1-t) + c2*t*t
    assert S.cancel(1 - 64*x*t*t*(1-t) - W/c0*(1-(1-u)*t)) == 0

    P = (
        c0*(k+S.Rational(1, 2))*(k+S.Rational(3, 2))
        + c1*(2*k+S.Rational(1, 2))*(k+S.Rational(1, 2))
        + c2*(2*k+S.Rational(1, 2))*(2*k+S.Rational(3, 2))
    )
    R = c0*(-2*k-2) - c1*k + c2*(4*k+2)
    assert S.expand(P - target_pr[0]) == 0
    assert S.expand(R - target_pr[1]) == 0

    U = S.cancel(a*DDB + b*DB + c*B)
    V = S.cancel(-4*a*DB*L - 2*a*B*DL - 2*b*B*L + d*DB + e*B)
    assert S.rem(S.fraction(U)[0], K, w) == 0
    assert S.rem(S.fraction(V)[0], K, w) == 0
    assert S.simplify(U.subs(w, S.sqrt(u)) - expected_u) == 0
    assert S.simplify(V.subs(w, S.sqrt(u)) - expected_v) == 0
    assert S.rem((1-w*w)**3 + 64*x*w*w, K, w) == 0
    assert S.simplify(((1-w*w)**3 + 64*x*w*w).subs(w, S.sqrt(u))) == 0
    print(f"PASS x={x}: convergence radius, W/Q, P/R, K cancellation, real-root values")

# Parameter/root transformations, checked independently of any special value.
xx = -(1-w*w)**3 / (64*w*w)
assert S.cancel(xx / S.diff(xx, w) + w*(1-w*w)/(2*(1+2*w*w))) == 0
r = S.symbols("r")
rr = 1 / (1-w*w)
assert S.cancel(((1-r)/(2-3*r)).subs(r, rr)/w - B) == 0
print("PASS universal root coefficient and Euler-operator transformations")
print("All independent exact algebraic checks passed.")
\end{filecontents*}
\begin{document}
\maketitle
\begin{abstract}
We give complete proofs of four numbered conjectures in Zhi-Wei Sun's
\emph{New series involving binomial coefficients (II)}, using the numbering
of arXiv:2307.03086v10: Conjectures 4.5, 4.7, 4.9, and 4.11.
Each conjecture consists of two identities, and all eight identities are
proved. Two continuous-parameter identities contain the eight evaluations
as specializations. A normalized beta integral converts parameter
derivatives into the combination $2H_{4k}-3H_{2k}+H_k$.
One family follows from exact cancellation of a quadratic factor;
the other follows from an algebraic differential operator that cancels
the contributions of two nonreal roots. All convergence arguments,
logarithm branches, polynomial factors, and remaining root values are
included. The source is standalone and embeds supplementary exact-check
programs. Historical priority and external refereeing are not claimed.
\end{abstract}

\section*{Scope and relation to the earlier work}
\addcontentsline{toc}{section}{Scope and relation to the earlier work}
The target numbering is fixed to Sun's revision of 6 February 2026
\cite{Sun}. Conjecture numbers and equation numbers are different.
\begin{center}
\begin{tabular}{@{}lll@{}}\toprule
Original conjecture & Original equations & Scope established here\\\midrule
4.5 &(4.9), (4.10) &Both identities\\
4.7 &(4.11), (4.12) &Both identities\\
4.9 &(4.13), (4.14) &Both identities\\
4.11 &(4.15), (4.16) &Both identities\\\bottomrule
\end{tabular}
\end{center}
These four source conjectures contain no further prime-congruence clauses.
In particular, the reference to source equation (4.13) does not mean
Conjecture 4.13.

The earlier research manuscript addressed Conjectures 4.19, 4.21, 4.25,
4.27, and 4.29. None of those five is counted again here. The transfer of
method is explicit: retain a free beta parameter, differentiate only after
establishing normal convergence, and finish with printed exact
certificates. The wider corpus was consulted for the practice of
constructing and independently verifying rational differential
certificates; no unproved assertion from that corpus is assumed.

The mathematical proofs below are complete arguments within ordinary
analysis and algebra. Independent internal checks and symbolic arithmetic
support the displayed certificates. They do not constitute external
peer review. A targeted literature comparison is recorded at the end,
with the limits of any priority inference stated separately.
\tableofcontents
\clearpage
\section{Four complete harmonic--binomial conjectures from one beta integral}
\label{sec:four-beta-closures}

\paragraph{Scope and provenance.}
This section proves both identities in each of Conjectures 4.5, 4.7, 4.9,
and 4.11 of Z.-W. Sun, \emph{New series involving binomial coefficients (II)},
\texttt{arXiv:2307.03086v10}, Section 4.
The corresponding equation numbers are (4.9)--(4.16).
These four numbered conjectures contain only these eight analytic identities;
no congruence assertion is being omitted.
The argument is self-contained. Its correctness and the historical priority
of the result are separate questions; no priority claim is made here.
Throughout,
\[
 H_0=0,\qquad H_n=\sum_{j=1}^n\frac1j,
 \qquad h_k=2H_{4k}-3H_{2k}+H_k,\qquad b_k=\binom{4k}{k}.
\]

\subsection{The eight identities}

\paragraph{Conjecture 4.5.}
\begin{align}
 \sum_{k=0}^{\infty}
 \frac{b_k\bigl((7k^2+10k+3)h_k-2k-4\bigr)}
 {(3k+1)(3k+2)24^k}
 &=-\sqrt3\log3, \label{eq:beta-sun49}\\
 \sum_{k=0}^{\infty}\frac{b_k}{24^k}
 \bigl((49k^2-146k+21)h_k-3038k+1160\bigr)
 &=\sqrt3(216-5\log3). \label{eq:beta-sun410}
\end{align}
\paragraph{Conjecture 4.7.}
\begin{align}
 \sum_{k=0}^{\infty}
 \frac{b_k\bigl((133k^2+131k+26)h_k+76k+40\bigr)}
 {(3k+1)(3k+2)(-25)^k}
 &=5\sqrt5\log5, \label{eq:beta-sun411}\\
 \sum_{k=0}^{\infty}\frac{b_k}{(-25)^k}
 \left((21413k^2-1409k+1036)h_k+
 \frac4{23}(118237k+17320)\right)
 &=\sqrt5\left(\frac{1440}{23}-100\log5\right).
 \label{eq:beta-sun412}
\end{align}
\paragraph{Conjecture 4.9.}
\begin{align}
 \sum_{k=0}^{\infty}
 \frac{b_k\bigl((55k^2+54k+11)h_k+22k+12\bigr)}
 {(3k+1)(3k+2)(-72)^k}
 &=3\sqrt3\log3, \label{eq:beta-sun413}\\
 \sum_{k=0}^{\infty}\frac{b_k}{(-72)^k}
 \left((3575k^2-1026k+67)h_k+
 \frac{242}{13}(175k+12)\right)
 &=\sqrt3\left(\frac{216}{13}-15\log3\right).
 \label{eq:beta-sun414}
\end{align}
\paragraph{Conjecture 4.11.}
\begin{align}
 \sum_{k=0}^{\infty}\frac{b_k}{128^k}
 \bigl((200k^2+76k-17)h_k-8(725k-49)\bigr)
 &=\sqrt2(144+5\log2), \label{eq:beta-sun415}\\
 \sum_{k=0}^{\infty}
 \frac{b_k\bigl((40k^2+44k+11)h_k-8(k+1)\bigr)}
 {(3k+1)(3k+2)128^k}
 &=-4\sqrt2\log2. \label{eq:beta-sun416}
\end{align}

\subsection{A parameter identity with its analytic domain}
Write $(a)_n=a(a+1)\cdots(a+n-1)$, with $(a)_0=1$. For
$0<\operatorname{Re}a<1$ and $|x|<27/256$, define
\begin{equation}
 F(a,x)=\sum_{k=0}^{\infty}
 \frac{64^k(a)_{2k}(1-a)_k}{(3k)!}x^k.
 \label{eq:beta-F}
\end{equation}
Euler's beta integral and the reflection identity for the gamma function give
\begin{equation}
 F(a,x)=\frac{\sin\pi a}{\pi}\int_0^1
 \frac{t^{a-1}(1-t)^{-a}}{Q_x(t)}\,dt,
 \qquad Q_x(t)=1-64xt^2(1-t).
 \label{eq:beta-euler}
\end{equation}
Indeed, the coefficient of $x^k$ on the right is
\[
 \frac{\sin\pi a}{\pi}64^k B(2k+a,k+1-a)
 =\frac{64^k(a)_{2k}(1-a)_k}{(3k)!}.
\]
For justification, $\max_{0\le t\le1}t^2(1-t)=4/27$, so the
geometric expansion is uniformly convergent whenever $|x|<27/256$.
On compact subsets of $0<\operatorname{Re}a<1$, its derivatives in $a$
are dominated by an integrable endpoint weight times a power of
$|\log t|+|\log(1-t)|$. Differentiation in $a$ and any fixed number of
applications of $D=x\,d/dx$ are therefore valid. In particular,
\begin{equation}
 F(1/2,x)=\sum_{k\ge0}b_kx^k,
 \qquad G(x):=\left.\partial_aF(a,x)\right|_{a=1/2}
 =\sum_{k\ge0}b_kh_kx^k.
 \label{eq:beta-G}
\end{equation}
The first equality follows from
$(1/2)_n=(2n)!/(4^n n!)$. For the second, logarithmic differentiation
of the Pochhammer factors gives
\[
 \sum_{j=0}^{2k-1}\frac1{j+1/2}-
 \sum_{j=0}^{k-1}\frac1{j+1/2}
 =2H_{4k}-3H_{2k}+H_k.
\]
These also prove absolute convergence of all eight series above.

\subsection{The four identities with quadratic denominators}
For constants $c_0,c_1,c_2$, put
\begin{equation}
 W(t)=c_0(1-t)^2+c_1t(1-t)+c_2t^2
 \quad\hbox{and}\quad
 J_W(a,x)=\frac{\sin\pi a}{\pi}
 \int_0^1\frac{t^{a-1}(1-t)^{-a}W(t)}{Q_x(t)}\,dt.
 \label{eq:beta-J}
\end{equation}
Termwise expansion of the three terms of $W$ gives
\begin{equation}
 J_W(a,x)=\sum_{k=0}^{\infty}\frac{64^kx^k}{(3k+2)!}
 \left[c_0(a)_{2k}(1-a)_{k+2}
 +c_1(a)_{2k+1}(1-a)_{k+1}
 +c_2(a)_{2k+2}(1-a)_k\right].
 \label{eq:beta-J-series}
\end{equation}
Consequently
\begin{equation}
 \left.\partial_aJ_W(a,x)\right|_{a=1/2}
 =\sum_{k=0}^{\infty}\frac{b_kx^k}{(3k+1)(3k+2)}
 \bigl(P_c(k)h_k+R_c(k)\bigr),
 \label{eq:beta-J-diff}
\end{equation}
where
\begin{align}
 P_c(k)&=c_0(k+\tfrac12)(k+\tfrac32)
 +c_1(2k+\tfrac12)(k+\tfrac12)
 +c_2(2k+\tfrac12)(2k+\tfrac32),\\
 R_c(k)&=c_0(-2k-2)-c_1k+c_2(4k+2).
\end{align}
These formulas include $k=0$: there is no omitted initial term.

\paragraph{A reusable one-parameter theorem.}
Let $u>0$ and put
\[
 x=\frac{(u-1)^3}{64u},\qquad |x|<\frac{27}{256},\qquad
 W_u(t)=1-(u-1)t+\frac{(u-1)^2}{u}t^2.
\]
Then
\[
 W_u(t)=(1-t)^2+(3-u)t(1-t)+u^{-1}t^2,
 \qquad Q_x(t)=(1+(u-1)t)W_u(t).
\]
Therefore, for every $a$ with $0<\operatorname{Re}a<1$,
\begin{equation}
 J_{W_u}(a,x)=u^{-a}.\label{eq:beta-all-u-parameter}
\end{equation}
Here cancellation is valid even at $u=1$, where $x=0$ and $W_u=1$.
In particular,
\begin{equation}
 \sum_{k=0}^\infty\frac{b_kx^k}
 {(3k+1)(3k+2)}\bigl(P_u(k)h_k+R_u(k)\bigr)
 =-\frac{\log u}{\sqrt u},\label{eq:beta-all-u-harmonic}
\end{equation}
where
\begin{align*}
 P_u(k)&=\left(7-2u+\frac4u\right)k^2
 +\left(\frac{13}2-\frac{3u}2+\frac4u\right)k
 +\frac32-\frac u4+\frac{3}{4u},\\
 R_u(k)&=\left(u-5+\frac4u\right)k-2+\frac2u.
\end{align*}
To prove \eqref{eq:beta-all-u-parameter}, cancel $W_u$ inside
\eqref{eq:beta-J} and set $y=t/(1-t)$, obtaining
\[
 \frac{\sin\pi a}{\pi}\int_0^\infty
 \frac{y^{a-1}}{1+uy}\,dy=u^{-a}.
\]
Then apply \eqref{eq:beta-J-diff} with
$(c_0,c_1,c_2)=(1,3-u,u^{-1})$.
The four concrete denominator identities below are scalar multiples
of this general theorem.

For the following four rows, direct polynomial multiplication gives
\begin{equation}
 Q_x(t)=\frac{W(t)}{c_0}\left(1-\frac t r\right),
 \qquad u=1-\frac1r>0.
 \label{eq:beta-cancel}
\end{equation}
\begin{center}
\begin{tabular}{c|c|c|c|c|c}
Sun equation &$x$&$(c_0,c_1,c_2)$&$W(t)$&$r$&$u$\\\hline
(4.9)&$1/24$&$(3,0,1)$&$4t^2-6t+3$&$-1/2$&$3$\\
(4.11)&$-1/25$&$(5,14,25)$&$16t^2+4t+5$&$5/4$&$1/5$\\
(4.13)&$-1/72$&$(3,8,9)$&$4t^2+2t+3$&$3/2$&$1/3$\\
(4.16)&$1/128$&$(8,8,4)$&$4(t^2-2t+2)$&$-1$&$2$
\end{tabular}
\end{center}
Thus the entire parameter identity reduces to
\begin{equation}
 J_W(a,x)=c_0\frac{\sin\pi a}{\pi}
 \int_0^1\frac{t^{a-1}(1-t)^{-a}}{1-t/r}\,dt
 =c_0\left(1-\frac1r\right)^{-a}=c_0u^{-a}.
 \label{eq:beta-parameter-collapse}
\end{equation}
For completeness, the last integral can be evaluated without
hypergeometric notation. Set $y=t/(1-t)$; the integrand becomes
$y^{a-1}/(1+uy)$ on $(0,\infty)$. Scaling $v=uy$ reduces it to
$u^{-a}B(a,1-a)$. Here $u>0$, so only the real logarithm is involved.
Differentiating at $a=1/2$ gives
\begin{equation}
 \sum_{k\ge0}\frac{b_kx^k(P_c(k)h_k+R_c(k))}
 {(3k+1)(3k+2)}=-\frac{c_0\log u}{\sqrt u}.
 \label{eq:beta-first-family-final}
\end{equation}
Substitution of the four rows gives exactly
\eqref{eq:beta-sun49}, \eqref{eq:beta-sun411},
\eqref{eq:beta-sun413}, and \eqref{eq:beta-sun416}.

\subsection{A finite algebraic formula for the remaining four sums}
Let $x\ne0$, $|x|<27/256$, and let $r_1,r_2,r_3$ be the three
roots of $Q_x$. They are distinct: a multiple root would require
$r=2/3$ and $x=27/256$, which is outside the open disk. None lies
on $[0,1]$, since $|64xt^2(1-t)|<1$ there. Partial fractions give
\begin{equation}
 \frac1{Q_x(t)}=\sum_{j=1}^3\frac{A(r_j)}{1-t/r_j},
 \qquad A(r)=\frac{1-r}{2-3r}.
 \label{eq:beta-partial}
\end{equation}
For each $r_j$, set
\[
 u_j=1-1/r_j,\qquad w_j=\sqrt{u_j},
\]
using the principal square root. Because $r_j\notin[0,1]$,
$u_j\notin(-\infty,0]$, so $\operatorname{Re}w_j>0$.
The same Euler integral used above, continued analytically to
$u\in\mathbb C\setminus(-\infty,0]$, yields
\begin{equation}
 F(a,x)=\sum_{j=1}^3 A(r_j)u_j^{-a}.
 \label{eq:beta-F-roots}
\end{equation}
This analytic continuation is legitimate because the integral
$\int_0^\infty y^{a-1}/(1+uy)\,dy$ and $u^{-a}B(a,1-a)$
are holomorphic on the same slit plane and agree for $u>0$.
All logarithms below are therefore fixed by these principal branches.

In the variable $w$, the root equation and the coefficient become
\begin{equation}
 x=-\frac{(1-w^2)^3}{64w^2},\qquad
 B(w):=\frac{A(r)}w=\frac{w}{1+2w^2}.
 \label{eq:beta-w}
\end{equation}
It follows that
\begin{equation}
 F(1/2,x)=\sum_{j=1}^3B(w_j),\qquad
 G(x)=-2\sum_{j=1}^3B(w_j)\log w_j.
 \label{eq:beta-root-sums}
\end{equation}
Locally at any nonzero $x$ with $|x|<27/256$, each root is analytic
and the differential operator $D=x\,d/dx$ becomes
\begin{equation}
 \mathcal D=-\frac{w(1-w^2)}{2(1+2w^2)}\frac d{dw},
 \qquad L(w):=\mathcal D\log w
 =-\frac{1-w^2}{2(1+2w^2)}.
 \label{eq:beta-D}
\end{equation}
There is no singularity in these formulas at the three selected roots:
$w=0$ cannot satisfy $(1-w^2)^3+64xw^2=0$, while $w^2=-1/2$
would force $x=27/256$, outside the assumed open disk.

For constants $a,b,c,d,e$, define
\begin{align}
 U(w)&=(a\mathcal D^2+b\mathcal D+c)B(w),\\
 V(w)&=-4a(\mathcal DB)L-2aB\mathcal DL-2bBL
       +d\mathcal DB+eB. \label{eq:beta-UV}
\end{align}
Then the product rule proves the exact finite formula
\begin{equation}
 \sum_{k=0}^{\infty}b_kx^k
 \bigl((ak^2+bk+c)h_k+dk+e\bigr)
 =\sum_{j=1}^3\bigl(-2U(w_j)\log w_j+V(w_j)\bigr).
 \label{eq:beta-polynomial-master}
\end{equation}
The constants for the four identities are
\begin{center}
\begin{tabular}{c|c|r|r|r|c|c|c}
Sun equation&$x$&$a$&$b$&$c$&$d$&$e$&$u$\\\hline
(4.10)&$1/24$&49&$-146$&21&$-3038$&1160&3\\
(4.12)&$-1/25$&21413&$-1409$&1036&$472948/23$&$69280/23$&$1/5$\\
(4.14)&$-1/72$&3575&$-1026$&67&$42350/13$&$2904/13$&$1/3$\\
(4.15)&$1/128$&200&76&$-17$&$-5800$&392&2
\end{tabular}
\end{center}

\subsection{A polynomial family containing the remaining four identities}
\label{sec:general-polynomial-family}

The specializations above arise from an identity with a free real parameter.
Suppose that $u>0$ and
\[
 x=x(u):=\frac{(u-1)^3}{64u},\qquad |x(u)|<\frac{27}{256}.
\]
Define the following polynomials in $u$:\nopagebreak
\begin{align*}
 A_u&=4(u-4)(2u-5)^2(2u+1)^2,\\*
 B_u&=2(2u-5)(12u^4-76u^3+237u^2-267u+40),\\*
 C_u&=(u-1)^2(2u-5)(4u^2-24u+15),\\*
 D_u&=-4(u-4)(u-1)(2u+25)(2u+1)^2,\\*
 E_u&=-8(u-1)^2(28u^2-44u+25).
\end{align*}
Then
\begin{equation}
\begin{split}
 &\sum_{k=0}^{\infty}b_k x(u)^k
 \left((A_uk^2+B_uk+C_u)h_k+D_uk+E_u\right)\\
 &\hspace{15mm}=\sqrt{u}(u-1)^2
 \left(5(2u-5)\log u-72u\right).
\end{split}
\label{eq:general-polynomial-family}
\end{equation}
At $u=1$ the identity is read directly at $x=0$ and both sides vanish.
There is no excluded value at $u=5/2$: the displayed polynomial coefficients
make that specialization meaningful as well.

Here is an explicit algebraic proof. Use the notation $s=w^2$ and set
\[
 K_u(s)=us^2+u(u-3)s+1.
\]
The root equation factors as
\[
 (1-w^2)^3+64x(u)w^2=-\frac{s-u}{u}K_u(s).
\]
For $u\ne1$ in the stated domain one has $u<4$, because the function
$(u-1)^3/u$ is increasing for $u>1$ and equals $27/4$ at $u=4$.
The discriminant of $K_u$ is
\[
 u^2(u-3)^2-4u=u(u-1)^2(u-4)<0.
\]
Thus the three selected squared roots are $u$ and the two nonreal roots of
$K_u$, and the real selected root is $w=\sqrt u$.

In the definitions of $U,V$ in \eqref{eq:beta-UV}, take
$(a,b,c,d,e)=(A_u,B_u,C_u,D_u,E_u)$. Direct differentiation gives the exact
rational-function identities
\begin{align}
 U(w)&=\frac{wK_u(s)\mathcal R_u(s)}{(1+2s)^5},
 \label{eq:general-U-certificate}\\
 V(w)&=\frac{wK_u(s)\mathcal T_u(s)}{(1+2s)^4},
 \label{eq:general-V-certificate}
\end{align}
where
\begin{align*}
 \mathcal R_u(s)=-(2u-5)\bigl[&4(20u^2+176u-385)s^2\\
                   &-16(34u^2-95u-20)s
                    +5(28u^2-56u+1)\bigr],\\
 \mathcal T_u(s)=-72\bigl[&(8u^3+20u^2-58u+75)s
                     -2u(4u^2-32u+37)\bigr].
\end{align*}
These are polynomial identities after multiplication by the printed
denominators, so they also hold at $u=5/2$. They show that both nonreal-root
contributions to \eqref{eq:beta-polynomial-master} are zero. At the remaining
root, elementary substitution gives
\[
 U(\sqrt u)=-5\sqrt u\,(u-1)^2(2u-5),\qquad
 V(\sqrt u)=-72u\sqrt u\,(u-1)^2.
\]
The value $-2U(\sqrt u)\log\sqrt u+V(\sqrt u)$ is precisely the right-hand
side of \eqref{eq:general-polynomial-family}. This proves the identity for
$u\ne1$; its $u=1$ case was already checked directly.

For the four original identities, the vectors
$(A_u,B_u,C_u,D_u,E_u)$ are scalar multiples of the coefficient rows in the
preceding table:
\[
 \begin{array}{c|c|c|c}
 \text{Sun equation}&u&x(u)&\text{scalar multiple}\\\hline
 (4.10)&3&1/24&-4\\
 (4.12)&1/5&-1/25&-92/3125\\
 (4.14)&1/3&-1/72&-52/243\\
 (4.15)&2&1/128&-1
 \end{array}
\]
Dividing by these nonzero scalars gives the four stated polynomial-weighted
identities. Together with the denominator family, this proves both assertions
in each of Sun's four numbered conjectures.

\subsection{Explicit factor certificates and the conclusion}
Put $s=w^2$ and write the common root polynomial as
\[
 E_x(w)=(1-w^2)^3+64xw^2.
\]
For each row there is a factorization
\begin{equation}
 E_x(w)=-\frac1\lambda(s-u)K(s).
 \label{eq:beta-E-factor}
\end{equation}
Here $(u,\lambda,K)$ are
\[
 \begin{array}{c|c|c|l}
 x&u&\lambda&K(s)\\\hline
 1/24&3&3&3s^2+1\\
 -1/25&1/5&5&5s^2-14s+25\\
 -1/72&1/3&3&3s^2-8s+9\\
 1/128&2&2&2s^2-2s+1.
 \end{array}
\]
Each displayed $K$ has negative discriminant. Thus the three values
$w_j^2$ consist of the positive real value $u$ and the two roots of $K$;
the selected real square root is $w_0=\sqrt u$.
Direct rational differentiation in \eqref{eq:beta-UV} gives the following
identities, valid as rational functions of $w$ (not merely at sampled
roots):
\begin{align*}
 x=1/24:\quad
 U(w)&=\frac{w(3s^2+1)(1292s^2-16s+425)}{4(1+2s)^5},\\
 V(w)&=\frac{54w(99s+46)(3s^2+1)}{(1+2s)^4};\\[2pt]
 x=-1/25:\quad
 U(w)&=\frac{w(5s^2-14s+25)(34900s^2-15056s+1135)}
 {4(1+2s)^5},\\
 V(w)&=\frac{18w(8033s-1538)(5s^2-14s+25)}{23(1+2s)^4};\\[2pt]
 x=-1/72:\quad
 U(w)&=\frac{w(3s^2-8s+9)(11668s^2-6896s+655)}
 {4(1+2s)^5},\\
 V(w)&=\frac{18w(1571s-482)(3s^2-8s+9)}{13(1+2s)^4};\\[2pt]
 x=1/128:\quad
 U(w)&=-\frac{w(2s^2-2s+1)(188s^2+1184s+5)}{(1+2s)^5},\\
 V(w)&=\frac{72w(103s+44)(2s^2-2s+1)}{(1+2s)^4}.
\end{align*}
Both $U$ and $V$ vanish at the two nonreal-root contributions.
The values at the remaining root are
\[
 \begin{array}{c|c|c|c}
 x&w_0&U(w_0)&V(w_0)\\\hline
 1/24&\sqrt3&5\sqrt3&216\sqrt3\\
 -1/25&1/\sqrt5&-100\sqrt5&1440\sqrt5/23\\
 -1/72&1/\sqrt3&-15\sqrt3&216\sqrt3/13\\
 1/128&\sqrt2&-5\sqrt2&144\sqrt2.
 \end{array}
\]
Formula \eqref{eq:beta-polynomial-master} therefore equals
\[
 -2U(\sqrt u)\log\sqrt u+V(\sqrt u)
 =-U(\sqrt u)\log u+V(\sqrt u).
\]
The four rows give exactly \eqref{eq:beta-sun410},
\eqref{eq:beta-sun412}, \eqref{eq:beta-sun414}, and
\eqref{eq:beta-sun415}. Together with the preceding beta-integral
collapse, this proves all eight identities and all four numbered
conjectures stated in this section.\hfill$\square$

\paragraph{Transfer from the earlier work.}
The reusable mechanism is the beta integral and differentiation with
respect to a parameter: the latter inserts the exact harmonic combination
$2H_{4k}-3H_{2k}+H_k$. Here the resulting denominator is cubic, and a
quadratic numerator cancels the two unwanted roots. This reduces the
integral to an elementary power $u^{-a}$. For the polynomially weighted
identities, the same cancellation takes place after applying the Euler
differential operator. This proof does not require a dilogarithm identity,
PSLQ, a truncation argument, or any of the eight asserted evaluations as
an input.

\section{Verification, literature comparison, and reproducibility}

\subsection{What was checked}
The coefficient normalization of the beta integral and its harmonic
derivative are proved for every nonnegative integer $k$. The geometric
majorant proves the required exchanges for the entire infinite series,
not only for a finite set of terms. The principal logarithms in the root
formula are justified on an explicitly specified slit plane.

The following algebraic assertions were also reconstructed and checked
independently over rational function fields: the quadratic cancellation
for each source specialization; the Euler differential operator in the
root variable; the two general-parameter factor certificates; their
real-root values; and the four scalar specializations. In particular,
the general certificates were checked with a symbolic parameter $u$,
not by substituting a finite collection of values.

As a separate transcription safeguard, all eight original sums were
evaluated using 400 terms at 100-decimal precision. Each discrepancy
from the printed value was below $1.9\times10^{-97}$. These numerical
comparisons are not used as a proof of any series identity.

\subsection{Literature comparison and the scope of the claims}
The targeted comparison was made on 19 November 2026 (Korea Standard
Time). It included the source conjecture paper \cite{Sun}, Au's WZ
papers \cite{Au2023,Au2024,Au2026}, Le Lan's beta-integral paper
\cite{LeLan}, Hou--Mu's WZ construction paper \cite{HouMu}, and
Hou--Sun's harmonic-series paper \cite{HouSun}.
The formulas were compared by their actual coefficients and harmonic
terms, in addition to their conjecture numbers.

Au's 2023 and 2024 arXiv revisions precede the February 2025 dates of
these four conjectures. His November 2026 harmonic-series paper was
also inspected; references there to similarly numbered identities in
Sun's other papers must not be confused with the targets here.
Le Lan proves source equation (4.19), belonging to Conjecture 4.15(ii),
and source equation (4.20), belonging to Conjecture 4.17(i), among
other evaluations. Those are different reciprocal-binomial sums;
neither numbered source conjecture is reduced to its analytic clause
when counting complete results.

A separate search target, Sun's Conjecture 5.31, was excluded from
the additional-result count after an exact prior proof was identified:
Hou--Sun's Theorem 1.23 proves all three of its identities, including
the evaluation $1316\zeta(5)$. Their Theorem 1.22 also proves all of
Conjecture 5.29. These known results are not claimed as new closures
in this manuscript. Neither theorem is one of the four targets proved
above.

No exact prior proof of the eight present target expressions was found
in that targeted check. This is not an exhaustive priority determination,
and it does not establish that the results remained open worldwide on
the date of this manuscript. The claim made here is the explicit,
self-contained proof and parameter extension of the four exact source
conjectures.

\subsection{Standalone source and supplementary programs}
All mathematics needed to read the proofs is printed in this PDF.
The accompanying single \texttt{.tex} file contains the entire document
and requires no external source file, image, bibliography database,
or data file. It can be compiled with a standard pdfLaTeX installation
and the packages listed in its preamble.

For reproducibility, the source embeds three supplementary Python
programs through \texttt{filecontents*}. Compiling the source writes
the following files into the compilation directory:
\begin{itemize}
\item \texttt{verify\_beta\_specializations.py};
\item \texttt{verify\_beta\_general\_families.py};
\item \texttt{verify\_beta\_independent.py}.
\end{itemize}
The programs require SymPy, and the first also requires mpmath. They
are optional checks; compilation does not execute them. The proofs
do not require trust in a numerical summation or a computer-algebra
black box.

\begin{thebibliography}{9}
\bibitem{Sun}
Z.-W. Sun, \emph{New series involving binomial coefficients (II)},
arXiv:2307.03086v10, 6 February 2026, Section 4.
\url{https://arxiv.org/html/2307.03086v10}.

\bibitem{Au2023}
K. C. Au, \emph{Multiple zeta values, WZ-pairs and infinite sums computations},
arXiv:2212.02986v4, 4 October 2023;
The Ramanujan Journal \textbf{66} (2025), article 3.
\url{https://arxiv.org/abs/2212.02986}.

\bibitem{Au2024}
K. C. Au, \emph{Wilf--Zeilberger seeds and non-trivial hypergeometric identities},
arXiv:2312.14051v4, 10 October 2024;
Journal of Symbolic Computation \textbf{130} (2025), 102421.
\url{https://arxiv.org/abs/2312.14051}.

\bibitem{Au2026}
K. C. Au, \emph{Discovering hypergeometric series with harmonic numbers
via Wilf--Zeilberger seeds}, arXiv:2602.08721v4, 9 November 2026.
\url{https://arxiv.org/abs/2602.08721}.

\bibitem{LeLan}
R. Le Lan, \emph{Evaluations of Series Using the Beta Function},
arXiv:2609.22304v2, 30 October 2026.
\url{https://arxiv.org/pdf/2609.22304v2}.

\bibitem{HouMu}
Q.-H. Hou and Y.-P. Mu, \emph{A construction method for WZ seeds},
arXiv:2604.22377, 24 April 2026.
\url{https://arxiv.org/abs/2604.22377}.

\bibitem{HouSun}
Q.-H. Hou and Z.-W. Sun, \emph{Taylor coefficients and series involving
harmonic numbers}, arXiv:2310.03699v3, 10 February 2026,
Theorems 1.22--1.23.
\url{https://arxiv.org/html/2310.03699v3}.
\end{thebibliography}
\end{document}
